\section{Introduction}

According to \cite{Rebelo:2005aa}, Finn Kydland and Edward Prescott's influential 1982 paper, ``Time to Build and Aggregate Fluctuations'', introduced three major ideas into macroeconomics. First, dynamic general equilibrium models provide a useful framework for studying business cycles. Second, business cycle models can be consistent with the empirical regularities of long-run growth. Third, calibrating models with microeconomic data and the economy's long-run properties allows researchers to compare simulated and observed data quantitatively.

David Backus, Patrick Kehoe and Finn Kydland applied these ideas in ``International Business Cycles: Theory and Evidence'' (1993),\footnote{The paper was later published in \textit{Frontiers of Business Cycle Research}, edited by Thomas F. Cooley (\cite{:1995aa}).} hereafter BKK or BKK (1993). They extended \cites{Kydland:1982aa} real business cycle (RBC) model to include a second country, calibrated and simulated the model, and compared its predictions with empirical regularities observed in the second half of the twentieth century. One discrepancy stood out for its robustness: cross-country output correlations were higher in the data than in the model, while consumption correlations were lower. Because this discrepancy persisted across calibrations and substantial changes to the model's structure, it became known as the quantity anomaly. These empirical findings shaped subsequent research, and many later dynamic stochastic general equilibrium (DSGE) models sought to account for them.

Given the importance of BKK's findings for research on international business cycles, this paper reassesses their empirical results using 23 additional years of data. \Cref{sec:themodel} introduces BKK's baseline model and compares it with its closed-economy counterpart. It also describes modifications that resolve some discrepancies between the model and the data but leave the quantity anomaly unresolved. \Cref{sec:thedata} describes our data, explains how they compare with BKK's, and discusses the implications of using the Hodrick-Prescott filter. In \cref{sec:theresults}, we repeat BKK's calculations over a longer sample period. We then examine all country pairs rather than only those involving the United States, comparing our results with both BKK's and those of \cite{Ambler:2004aa}. We find that the quantity anomaly persists and conclude by discussing how subsequent research has sought to explain it.

This paper also aims to make its analysis reproducible.\footnote{Replication is a central standard for assessing scientific research, but independently repeating every stage of an empirical study is not always feasible. Reproducible research provides an intermediate standard: publishing the data and code allows others to verify how the reported results were obtained. See \cite{Peng:2011aa}.} The data and the R code used in the original analysis are available online, together with instructions for reproducing the results.\footnote{The code, instructions and data are available at \url{https://denismaciel.github.io/BKK-International-Business-Cycle/}.}